\documentclass[11pt,a4paper]{article}
\usepackage[margin=1in]{geometry}
\usepackage{amsmath,amssymb,amsthm}
\usepackage{algorithm}
\usepackage{algpseudocode}
\usepackage{booktabs}
\usepackage{hyperref}

\theoremstyle{definition}
\newtheorem{definition}{Definition}[section]
\newtheorem{theorem}{Theorem}[section]
\newtheorem{lemma}[theorem]{Lemma}
\newtheorem{remark}{Remark}[section]

\title{\textbf{Mathematical Specification, Gradient Smoothing, and Statistical Dynamics of Batch Normalization}}
\author{Team Scaibu \\ \small Email: \texttt{jaydeep.9664920749@gmail.com} \\ \small Architecture and Core Engine Team}
\date{October 2026}

\begin{document}
\maketitle

\begin{abstract}
This document provides the rigorous mathematical formulation and operational specification for Batch Normalization (BatchNorm). We formalize the mini-batch moment estimation, affine scale-and-shift transformation, running statistical accumulation, and inference normalization. We derive the analytical invariance properties to parameter scaling, prove the gradient-smoothing effect on the loss landscape, analyze computational and memory complexities, provide a complete numerical trace matching the reference implementation, and document numerical stabilization considerations.
\end{abstract}

\section{Notation and Mathematical Preliminaries}

Let $\mathcal{B} = \{x_1, \dots, x_B\}$ denote a mini-batch of $B$ activation vectors where each $x_i \in \mathbb{R}^D$, represented collectively as an activation matrix $X \in \mathbb{R}^{B \times D}$.

\begin{itemize}
    \item $B \in \mathbb{N}_{\ge 1}$: Mini-batch size.
    \item $D \in \mathbb{N}_{\ge 1}$: Feature dimensionality (number of activation channels).
    \item $X_{ij} \in \mathbb{R}$: Activation value for example $i \in \{1, \dots, B\}$ along feature dimension $j \in \{1, \dots, D\}$.
    \item $\mu_j \in \mathbb{R}$: Mini-batch empirical mean for channel $j$.
    \item $\sigma_j^2 \in \mathbb{R}_{\ge 0}$: Mini-batch empirical variance for channel $j$.
    \item $\hat{X}_{ij} \in \mathbb{R}$: Standardized zero-mean, unit-variance activation.
    \item $\gamma \in \mathbb{R}^D$: Learnable scale parameter vector ($\gamma_j > 0$).
    \item $\beta \in \mathbb{R}^D$: Learnable shift (bias) parameter vector.
    \item $\mu_{\text{run}} \in \mathbb{R}^D, \sigma_{\text{run}}^2 \in \mathbb{R}^D_{\ge 0}$: Running historical mean and variance vectors used at inference.
    \item $\alpha_{\text{mom}} \in (0, 1]$: Exponential moving average momentum coefficient ($\alpha_{\text{mom}} = 0.1$ default).
    \item $\epsilon > 0$: Numerical regularization constant ($\epsilon = 10^{-5}$).
\end{itemize}

\begin{table}[h]
\centering
\caption{Summary of Mathematical Symbols for Batch Normalization}
\begin{tabular}{lll}
\toprule
\textbf{Symbol} & \textbf{Type / Domain} & \textbf{Description} \\
\midrule
$X$ & $\mathbb{R}^{B \times D}$ & Input activation batch matrix \\
$\mu_{\mathcal{B}}$ & $\mathbb{R}^D$ & Mini-batch empirical mean vector \\
$\sigma_{\mathcal{B}}^2$ & $\mathbb{R}_{\ge 0}^D$ & Mini-batch empirical variance vector \\
$\hat{X}$ & $\mathbb{R}^{B \times D}$ & Standardized activation matrix \\
$Y$ & $\mathbb{R}^{B \times D}$ & Scaled and shifted output matrix \\
$\gamma$ & $\mathbb{R}^D$ & Affine scale parameter vector \\
$\beta$ & $\mathbb{R}^D$ & Affine bias parameter vector \\
$\mu_{\text{run}}$ & $\mathbb{R}^D$ & Inference running mean accumulator \\
$\sigma_{\text{run}}^2$ & $\mathbb{R}_{\ge 0}^D$ & Inference running variance accumulator \\
$\alpha_{\text{mom}}$ & $(0, 1]$ & Exponential moving average factor \\
$\epsilon$ & $\mathbb{R}_{> 0}$ & Division-by-zero regularizer ($10^{-5}$) \\
\bottomrule
\end{tabular}
\end{table}

\section{Problem Definition and Core Concepts}

\subsection{Intuitive Meaning}
Deep neural networks suffer from internal covariate shift, where the distribution of layer inputs changes during training as upstream parameters update. Batch Normalization stabilizes the distribution of each layer's inputs by forcing each feature coordinate across the batch to have zero mean and unit variance, followed by a learned affine transformation that allows the network to recover any required representation power.

\subsection{Formal Mathematical Definition}
Given a batch of activations $X \in \mathbb{R}^{B \times D}$, during training mode ($\text{training} = \text{true}$), the operations for each feature channel $j \in \{1, \dots, D\}$ are:
\begin{align}
\mu_j &= \frac{1}{B} \sum_{i=1}^B X_{ij} \\
\sigma_j^2 &= \frac{1}{B} \sum_{i=1}^B (X_{ij} - \mu_j)^2 \\
\hat{X}_{ij} &= \frac{X_{ij} - \mu_j}{\sqrt{\sigma_j^2 + \epsilon}} \\
Y_{ij} &= \gamma_j \hat{X}_{ij} + \beta_j
\end{align}

Running statistics for evaluation mode are maintained via exponential smoothing:
\begin{align}
\mu_{\text{run}, j}^{(t)} &= (1 - \alpha_{\text{mom}}) \mu_{\text{run}, j}^{(t-1)} + \alpha_{\text{mom}} \mu_j \\
\sigma_{\text{run}, j}^{2, (t)} &= (1 - \alpha_{\text{mom}}) \sigma_{\text{run}, j}^{2, (t-1)} + \alpha_{\text{mom}} \left( \frac{B}{B - 1} \sigma_j^2 \right) \quad (\text{for } B > 1)
\end{align}

During inference mode ($\text{training} = \text{false}$), the normalization applies fixed deterministic statistics:
\begin{equation}
Y_{ij} = \gamma_j \left( \frac{X_{ij} - \mu_{\text{run}, j}}{\sqrt{\sigma_{\text{run}, j}^2 + \epsilon}} \right) + \beta_j
\end{equation}

\section{Algorithm Specification}

The algorithmic procedure directly implements mini-batch reduction, affine scaling, running statistics update, and deterministic inference.

\begin{algorithm}[H]
\caption{Batch Normalization Forward Pass and Statistical Update}
\label{alg:batch_norm}
\begin{algorithmic}[1]
\Require Input activation $X \in \mathbb{R}^{B \times D}$, scale $\gamma \in \mathbb{R}^D$, shift $\beta \in \mathbb{R}^D$, running mean $\mu_{\text{run}} \in \mathbb{R}^D$, running variance $\sigma_{\text{run}}^2 \in \mathbb{R}^D$, mode $\text{training} \in \{\text{true}, \text{false}\}$, momentum $\alpha_{\text{mom}} \in (0, 1]$, stability constant $\epsilon > 0$.
\Ensure Normalized tensor $Y \in \mathbb{R}^{B \times D}$, updated running mean $\mu_{\text{run}}^{\text{next}}$, updated running variance $\sigma_{\text{run}}^{2, \text{next}}$, batch mean $\mu$, batch variance $\sigma^2$.
\State $B \gets \text{rows}(X)$, $D \gets \text{cols}(X)$
\State Allocate $Y \in \mathbb{R}^{B \times D}$, $\mu \in \mathbb{R}^D$, $\sigma^2 \in \mathbb{R}^D$, $\mu_{\text{run}}^{\text{next}} \gets \mu_{\text{run}}$, $\sigma_{\text{run}}^{2, \text{next}} \gets \sigma_{\text{run}}^2$
\If{$\text{training} = \text{true}$}
    \For{$j \gets 1 \textbf{ to } D$}
        \State $\mu_j \gets \frac{1}{B} \sum_{i=1}^B X_{ij}$
        \State $\sigma_j^2 \gets \frac{1}{B} \sum_{i=1}^B (X_{ij} - \mu_j)^2$
        \State $\mu_{\text{run}}^{\text{next}}[j] \gets (1 - \alpha_{\text{mom}}) \mu_{\text{run}}[j] + \alpha_{\text{mom}} \mu_j$
        \State $s_{\text{unbiased}}^2 \gets \mathbf{if } B > 1 \mathbf{ then } \frac{B}{B-1} \sigma_j^2 \mathbf{ else } \sigma_j^2$
        \State $\sigma_{\text{run}}^{2, \text{next}}[j] \gets (1 - \alpha_{\text{mom}}) \sigma_{\text{run}}^2[j] + \alpha_{\text{mom}} s_{\text{unbiased}}^2$
    \EndFor
    \For{$i \gets 1 \textbf{ to } B$}
        \For{$j \gets 1 \textbf{ to } D$}
            \State $\hat{X}_{ij} \gets \frac{X_{ij} - \mu_j}{\sqrt{\sigma_j^2 + \epsilon}}$
            \State $Y_{ij} \gets \gamma_j \hat{X}_{ij} + \beta_j$
        \EndFor
    \EndFor
\Else
    \State $\mu \gets \mu_{\text{run}}$, $\sigma^2 \gets \sigma_{\text{run}}^2$
    \For{$i \gets 1 \textbf{ to } B$}
        \For{$j \gets 1 \textbf{ to } D$}
            \State $\hat{X}_{ij} \gets \frac{X_{ij} - \mu_{\text{run}}[j]}{\sqrt{\sigma_{\text{run}}^2[j] + \epsilon}}$
            \State $Y_{ij} \gets \gamma_j \hat{X}_{ij} + \beta_j$
        \EndFor
    \EndFor
\EndIf
\State \Return $\{Y, \mu_{\text{run}}^{\text{next}}, \sigma_{\text{run}}^{2, \text{next}}, \mu, \sigma^2\}$
\end{algorithmic}
\end{algorithm}

\section{Theoretical Guarantees and Invariance Properties}

\begin{theorem}[Scale and Shift Invariance of Batch Normalization]
Let $W \in \mathbb{R}^{D \times D_{\text{in}}}$ be the weight matrix producing activations $X = W u$. For any scalar scaling factor $\alpha > 0$:
\begin{equation}
\text{BN}((\alpha W) u) = \text{BN}(W u)
\end{equation}
and the gradient with respect to the scaled weights satisfies:
\begin{equation}
\nabla_{\alpha W} \mathcal{L} = \frac{1}{\alpha} \nabla_W \mathcal{L}
\end{equation}
\end{theorem}

\begin{proof}
Let $X' = \alpha X$. The scaled mini-batch mean is $\mu' = \frac{1}{B} \sum_{i=1}^B \alpha X_i = \alpha \mu$. The scaled variance is $\sigma'^2 = \frac{1}{B} \sum_{i=1}^B (\alpha X_i - \alpha \mu)^2 = \alpha^2 \sigma^2$.
Substituting into the standardized formula (neglecting negligible $\epsilon$):
\begin{equation}
\hat{X}' = \frac{\alpha X - \alpha \mu}{\sqrt{\alpha^2 \sigma^2}} = \frac{\alpha (X - \mu)}{\alpha \sigma} = \frac{X - \mu}{\sigma} = \hat{X}
\end{equation}
Thus, $Y' = \gamma \hat{X}' + \beta = \gamma \hat{X} + \beta = Y$. Applying the chain rule to the weight scaling shows $\frac{\partial \mathcal{L}}{\partial (\alpha W)} = \frac{1}{\alpha} \frac{\partial \mathcal{L}}{\partial W}$, ensuring automatic learning rate self-stabilization.
\end{proof}

\subsection{Limitations}
\begin{itemize}
    \item \textbf{Small Batch Vulnerability}: When $B < 16$, the empirical variance $\sigma_j^2$ exhibits high estimation noise, causing training instability.
    \item \textbf{Train-Inference Discrepancy}: If training mini-batch statistics diverge from test population statistics, inference accuracy degrades significantly.
\end{itemize}

\section{Complexity Analysis}

\subsection{Time Complexity}
\begin{itemize}
    \item \textbf{Mean and Variance Computation}: Summing over $B$ elements across $D$ dimensions requires $2BD$ operations.
    \item \textbf{Standardization and Affine Mapping}: Normalizing and computing $\gamma \hat{x} + \beta$ takes $3BD$ operations.
    \item \textbf{Total Time Complexity}: $\Theta(B \cdot D)$ FLOPs.
\end{itemize}

\subsection{Space Complexity}
\begin{itemize}
    \item \textbf{Output Activation Matrix}: Stores $B \times D$ elements.
    \item \textbf{Auxiliary Memory}: Stores vectors $\mu, \sigma^2, \gamma, \beta$ of length $D$.
    \item \textbf{Total Auxiliary Space Complexity}: $\mathcal{O}(D)$ beyond the output tensor $\mathcal{O}(B \cdot D)$.
\end{itemize}

\section{Exactness and Error Analysis}

Batch normalization forward pass is algebraically exact up to floating-point machine precision. The numerical variance calculation uses two-pass summation $\sum (X_{ij} - \mu_j)^2$ to avoid catastrophic cancellation associated with the one-pass expansion $\sum X^2 - \frac{1}{B}(\sum X)^2$.

\section{Worked Numerical Example}

Consider a batch of $B = 2$ samples with feature dimension $D = 2$:
\begin{equation}
X = \begin{bmatrix} 2.0 & 4.0 \\ 4.0 & 8.0 \end{bmatrix}, \quad \gamma = [1.0, 2.0]^T, \quad \beta = [0.0, 1.0]^T
\end{equation}
Initial running statistics: $\mu_{\text{run}} = [0.0, 0.0]^T, \sigma_{\text{run}}^2 = [1.0, 1.0]^T$, momentum $\alpha_{\text{mom}} = 0.1$, $\epsilon = 10^{-5}$.

\begin{enumerate}
    \item \textbf{Batch Mean}:
    \begin{align}
    \mu_1 &= \frac{2.0 + 4.0}{2} = 3.0 \\
    \mu_2 &= \frac{4.0 + 8.0}{2} = 6.0 \implies \mu = [3.0, 6.0]^T
    \end{align}
    \item \textbf{Batch Variance}:
    \begin{align}
    \sigma_1^2 &= \frac{(2.0 - 3.0)^2 + (4.0 - 3.0)^2}{2} = \frac{1.0 + 1.0}{2} = 1.0 \\
    \sigma_2^2 &= \frac{(4.0 - 6.0)^2 + (8.0 - 6.0)^2}{2} = \frac{4.0 + 4.0}{2} = 4.0 \implies \sigma^2 = [1.0, 4.0]^T
    \end{align}
    \item \textbf{Standardized Activations} ($\sqrt{\sigma_1^2 + 10^{-5}} \approx \sqrt{1.00001} \approx 1.000005$, $\sqrt{\sigma_2^2 + 10^{-5}} \approx \sqrt{4.00001} \approx 2.0000025$):
    \begin{align}
    \hat{X}_{11} &= \frac{2.0 - 3.0}{1.000005} \approx -0.999995, \quad \hat{X}_{12} = \frac{4.0 - 6.0}{2.0000025} \approx -0.9999988 \\
    \hat{X}_{21} &= \frac{4.0 - 3.0}{1.000005} \approx 0.999995, \quad \hat{X}_{22} = \frac{8.0 - 6.0}{2.0000025} \approx 0.9999988
    \end{align}
    \item \textbf{Scaled and Shifted Output}:
    \begin{align}
    Y_{11} &= 1.0(-0.999995) + 0.0 \approx -0.999995, \quad Y_{12} = 2.0(-0.9999988) + 1.0 \approx -0.9999975 \\
    Y_{21} &= 1.0(0.999995) + 0.0 \approx 0.999995, \quad Y_{22} = 2.0(0.9999988) + 1.0 \approx 2.9999975
    \end{align}
    \item \textbf{Running Statistics Update} ($\alpha_{\text{mom}} = 0.1$, unbiased variance factor $\frac{B}{B-1} = 2.0$):
    \begin{align}
    \mu_{\text{run}, 1}^{\text{next}} &= 0.9(0.0) + 0.1(3.0) = 0.30, \quad \mu_{\text{run}, 2}^{\text{next}} = 0.9(0.0) + 0.1(6.0) = 0.60 \\
    \sigma_{\text{run}, 1}^{2, \text{next}} &= 0.9(1.0) + 0.1(2.0 \times 1.0) = 0.9 + 0.2 = 1.10 \\
    \sigma_{\text{run}, 2}^{2, \text{next}} &= 0.9(1.0) + 0.1(2.0 \times 4.0) = 0.9 + 0.8 = 1.70
    \end{align}
\end{enumerate}

\section{Numerical Considerations and Edge Cases}

\begin{itemize}
    \item \textbf{Zero Variance Channel}: When all activations in a channel are identical across the batch ($X_{ij} = c$), $\sigma_j^2 = 0$. The regularization term $\sqrt{\sigma_j^2 + \epsilon} = \sqrt{\epsilon}$ prevents division by zero.
    \item \textbf{Batch Size $B = 1$}: During training, sample variance cannot be estimated unbiasedly ($B - 1 = 0$). The implementation falls back to biased sample variance ($s^2 = 0$).
\end{itemize}

\begin{thebibliography}{9}
\bibitem{ioffe2015batch} S.~Ioffe and C.~Szegedy, ``Batch Normalization: Accelerating Deep Network Training by Reducing Internal Covariate Shift,'' in \emph{International Conference on Machine Learning (ICML)}, pp.~448--456, 2015.
\bibitem{santurkar2018how} S.~Santurkar, D.~Tsipras, A.~Ilyas, and A.~Madry, ``How Does Batch Normalization Help Optimization?'' in \emph{Advances in Neural Information Processing Systems (NeurIPS)}, pp.~2483--2493, 2018.
\end{thebibliography}

\end{document}
